% !TEX root = ../main.tex

\chapter{Introdução}
\label{chap:introducao}

A Geometria Diferencial, área amplamente desenvolvida por Carl Friedrich Gauss, busca generalizar conceitos do cálculo diferencial e integral para superfícies, entendidas como subconjuntos do espaço real tridimensional que localmente se comportam como um espaço bidimensional. Um dos resultados mais importantes obtidos por Gauss é o Teorema Egrégio, segundo o qual a curvatura de uma superfície é uma propriedade intrínseca, isto é, independe do espaço ambiente e pode ser determinada exclusivamente a partir da métrica definida na própria superfície \cite{gauss1902}. Tal resultado representa um passo fundamental na direção da concepção de objetos geométricos independentes de um espaço ambiente e, por consequência, no desenvolvimento de uma teoria mais geral.

Em 1854, Bernhard Riemann propôs uma generalização dessas ideias ao questionar o porquê da restrição do estudo geométrico a espaços localmente semelhantes ao espaço real bidimensional. A partir dessa reflexão, introduziu a noção de espaços que localmente se comportam como $\mathbb{R}^{n}$, equipados com uma noção infinitesimal de distância \cite{riemann1873}. Décadas depois, com a contribuição de diversos matemáticos, o conceito de variedade diferenciável foi formalizado, permitindo o estudo rigoroso da geometria intrínseca de espaços de dimensões finitas arbitrárias, sem a necessidade de um espaço ambiente, e servindo de modelo para geometrias abstratas, em particular, as não-euclidianas \cite{lima1960variedades}.

Um conceito fundamental da Geometria Riemanniana é o de geodésica. De maneira intuitiva, geodésicas são curvas que generalizam a noção de reta para variedades diferenciáveis e que, localmente, realizam a menor distância entre dois pontos. Em outras palavras, é o objeto da Geometria Riemanniana que tem papel análogo ao segmento de reta na geometria clássica \cite{doCarmo1992}.

Esse trabalho desenvolverá toda a teoria necessária para um estudo aprofundado das geodésicas, incluindo diferentes formas de classificá-las e a construção de certas variedades nas quais as geodésicas estão definidas globalmente. O texto é auto-contido, com exceção de alguns teoremas clássicos da topologia, análise no $\mathbb{R}^{n}$, algebra linear e equações diferenciais, os quais estão propriamente referenciados.

Começamos, no \cref{cap:preliminares}, por apresentar alguns resultados da topologia diferencial, como a definição de variedade diferenciável e diferenciabilidade, além de demonstrar certas maneiras de construir variedades as quais serão úteis ao longo do texto. No \cref{cap:metricas-conexoes}, introduzimos a geometria riemanniana por apresentar a noção de variedade riemanniana, a qual é uma variedade diferenciável, onde todo plano tangente está munido com um produto interno que varia diferenciavelmente com respeito aos pontos da variedade. Também exibimos a noção de conexão, a qual serve para diferenciarmos campos de vetores ao longo de outros campos de vetores em uma variedade riemanniana, e é fundamental para definir e analisar as geodésicas. Finalizamos o capítulo demonstrando a existência e unicidade, em qualquer variedade riemanniana, de uma conexão simétrica e compatível com a métrica. Por fim, no \cref{cap:geodesicas-hopf-rinow}, definimos as geodésicas, mostramos que por qualquer ponto e com qualquer velocidade, passa por ele uma geodésica, com tal velocidade, definida em um intervalo de tempo suficientemente pequeno. Mostramos também que toda curva que realiza a menor distância entre dois pontos e é parametrizada por comprimento de arco deve ser uma geodésica, e a recíproca local desse resultado, que afirma que toda geodésica realiza a menor distância localmente. Concluímos por tornar da variedade um espaço métrico e demonstramos o \Cref{teo:hopf-rinow}, o qual afirma, dentre outras coisas, que a completude métrica é equivalente a completude geodésica e a mesma pode ser obtida analisando o mapa exponencial no plano tangente de um único ponto arbitrário.